%!TEX root = ../bachelorarbeit.tex

\addchap{Abkürzungsverzeichnis}
%nur verwendete Akronyme werden letztlich im Abkürzungsverzeichnis des Dokuments angezeigt
%Verwendung: 
%		\ac{Abk.}   --> fügt die Abkürzung ein, beim ersten Aufruf wird zusätzlich automatisch die ausgeschriebene Version davor eingefügt bzw. in einer Fußnote (hierfür muss in header.tex \usepackage[printonlyused,footnote]{acronym} stehen) dargestellt
%		\acs{Abk.}   -->  fügt die Abkürzung ein
%		\acf{Abk.}   --> fügt die Abkürzung UND die Erklärung ein
%		\acl{Abk.}   --> fügt nur die Erklärung ein
%		\acp{Abk.}  --> gibt Plural aus (angefügtes 's'); das zusätzliche 'p' funktioniert auch bei obigen Befehlen
%	siehe auch: http://golatex.de/wiki/%5Cacronym
%	
\begin{acronym}[YTMMM]

%To-Do: Reihenfolge der Abkürzungen anpassen, damit sie alphabetisch sortiert sind oder der Reihenfolge im Text entsprechen

\acro{KI}{Künstliche Intelligenz}
\acro{ESG}{Environmental, Social and Governance}
\acro{YoY}{Year-over-Year}
\acro{LLM}{Large Language Model}
\acro{DSR}{Design Science Research}
\acro{XAI}{Explainable Artificial Intelligence}
\acro{MPT}{Moderne Portfoliotheorie}
\acro{DRL}{Deep Reinforcement Learning}
\acro{ML}{Machine Learning}
\acro{EU}{Europäische Union}

\acro{SHAP}{SHapley Additive exPlanations}
\acro{LIME}{Local Interpretable Model-agnostic Explanations}
\acro{CoT}{Chain-of-Thought}
\acro{GPT}{Generative Pre-trained Transformer}
\acro{BERT}{Bidirectional Encoder Representations from Transformers}
\acro{SaaS}{Software as a Service}
\acro{ETF}{Exchange-Traded Fund}
\acro{DSGVO}{Datenschutzgrundsatzverordnung}

\acro{API}{Application Programming Interface}
\acro{CLI}{Command Line Interface}
\acro{CPU}{Central Processing Unit}
\acro{DB}{Datenbank}
\acro{DBMS}{Datenbankmanagementsystem}
\acro{DDL}{Data Definition Language}
\acro{DML}{Data Manipulation Language}
\acro{FOSS}{Free and Open Source Software}
\acro{GUI}{Graphical User Interface}
\acro{HTML}{Hypertext Markup Language}
\acro{HTTP}{Hypertext Transfer Protocol}
\acro{JDBC}{Java Database Connectivity}
\acro{JSON}{JavaScript Object Notation}
\acro{JS}{JavaScript}
\acro{JVM}{Java Virtual Machine}
\acro{PoC}{Proof of Concept}
\acro{RAM}{Random Access Memory}
\acro{REST}{Representational State Transfer}
\acro{PDF}{Portable Document Format}
\acro{TCP}{Transmission Control Protocol}
\acro{TXT}{Textdatei}
\acro{VCS}{Versionskontrollsystem}

\end{acronym}
